\documentclass{article}
\usepackage{graphicx} % Required for inserting images
\usepackage[russian]{babel}
\usepackage{cancel}
\usepackage{amsmath}
\usepackage{listings}
\usepackage{fancyhdr}
\usepackage{caption}
\usepackage{pdfpages}
\usepackage{amssymb}
\captionsetup[figure]{font=large}
\DeclareMathOperator{\erf}{erf}
\usepackage{setspace} % Подключаем пакет для настройки интервала
\onehalfspacing % Устанавливаем интервал 1.5
\usepackage[
    left=30mm,      % Левое поле 30 мм
    right=15mm,     % Правое поле 15 мм
    top=20mm,       % Верхнее поле 20 мм (опционально)
    bottom=20mm,    % Нижнее поле 20 мм (опционально)
]{geometry}

\begin{document}
\Large



\section{Постановка задачи}
\subsection{Прямая задача}
Рассмотрим задачу распространения тепла в полуограниченном стержне при наличии внутренних источников, плотность мощности которых описывается функцией $f(x, t)$. Начальное распределение температуры стержня задано и равно $u_0(x)$, а на границе $x = 0$ определен тепловой поток $\gamma(t)$. 
Математическая формулировка задачи имеет следующий вид:

\begin{align*}
        &u_t(x, t)= a^2u_{xx}(x, t) + f(x, t), \ \ x > 0, \ \ t > 0, \\
        &u_x(0, t) = \gamma(t), \ \ t > 0,\\
        &u(x, 0) = u_0(x), \ \ x >  0.
\end{align*}

\subsection{Обратная задача}
Рассмотрим теперь обратную задачу определения источника. Слагаемое $f(x,t) \equiv f(x)$ считается неизвестным. Для нахождения источника фиксируется дополнительная информация, а именно - значение функции $u_t(x, t)$ на правой границе отрезка времени $[0, T]$. Постановка задачи имеет вид.

\begin{align*}
        &u_t(x, t)= a^2u_{xx}(x, t) + f(x), \ \ x > 0, \ \ t > 0, \\
        &u_x(0, t) = \gamma(t), \ \ t > 0,\\
        &u(x, 0) = u_0(x), \ \ x > 0,\\
        &u_t(x, T) = u_T(x), \ \ x > 0.
\end{align*}

Функция $u_T(x)$ задана. Таким образом, зная $\gamma(t), u_0(x), u_T(x)$, требуется восстановить $u(x, t)$ и $f(x)$.
\section{Исследование задачи}
В силу линейности исходной задачи путём подстановок ее можно редуцировать к более простой.

\begin{align*}
        &u_t(x, t)= a^2u_{xx}(x, t) + f(x), \ \ x  > 0, \ \ t > 0, \\
        &u_x(0, t) = 0, \ \ t > 0,\\
        &u(x, 0) = 0, \ \ x > 0,\\
        &u_t(x, T) = \varphi(x), \ \ x > 0.
\end{align*}

\subsection{Эквивалентность уравнению теплопроводности с обратным направлением времени}
Заметим следующее. Если перевести задачу в Фурье образы
 $$ \Phi(\xi) =  \int_{0}^{+\infty} \varphi(x)\cos(\xi x) \,dx, $$
 $$ \text{F}(\xi) =  \int_{0}^{+\infty} f(x)\cos(\xi x) \,dx, $$
тогда уравнение будет иметь вид

\begin{align*}
    \Phi(\xi) = e^{-a^2\xi^2T}\text{F}(\xi).
\end{align*}

Продолжив функции исходной задачи четным образом на всю прямую $\mathbb{R}$, можно получить эквивалентную постановку.

\begin{align*}
        &u_t(x, t)= a^2u_{xx}(x, t), \ \ x  \in \mathbb{R}, \ \ t > 0, \\
        &u(-x, t) = u(x, t), \ \ x  \in \mathbb{R}, \ \  t > 0,\\
        &u(x, T) = \varphi(x), \ \ x  \in \mathbb{R}.
\end{align*}

Требуется восстановить $u_0(x) = u(x, 0), \ \ x \in \mathbb{R}$.
Тогда решением исходной обратной задачи будет $f(x) = u_0(x), \ \ x > 0$.
Далее будем рассматривать лишь эквивалетную постановку.
\subsection{Разрешающие формулы}
$$u(x, t) = \int_{0}^{t}\frac{1}{\sqrt{4\pi a^2 s}}\int_{0}^{+\infty}\bigg(e^{-\frac{(x - y)^2}{4a^2s}} + e^{-\frac{(x + y)^2}{4a^2s}}\bigg)f(y) \, dy \, ds.$$
Дифференцирую по переменной $t$, получим
$$u_t(x, t) = \frac{1}{\sqrt{4\pi a^2 t}}\int_{0}^{+\infty}\bigg(e^{-\frac{(x - y)^2}{4a^2t}} + e^{-\frac{(x + y)^2}{4a^2t}}\bigg)f(y) \, dy$$
тогда
$$\varphi(x) = u_t(x, T) = \frac{1}{\sqrt{4\pi a^2 T}}\int_{0}^{+\infty}\bigg(e^{-\frac{(x - y)^2}{4a^2T}} + e^{-\frac{(x + y)^2}{4a^2T}}\bigg)f(y) \, dy. = A_T(f)(x)$$
Итого для восстановления $f(x)$ требуется решить интегральное уравнение Фредгольма первого рода.

\end{document}